\documentclass[11pt,a4paper,oneside]{report}

% --- Packages ---
\usepackage[utf8]{inputenc}
\usepackage[T1]{fontenc}
\usepackage{lmodern}
\usepackage[margin=2.2cm,headheight=14pt]{geometry}
\usepackage{xcolor}
\usepackage{graphicx}
\usepackage{enumitem}
\usepackage{parskip}
\usepackage{titlesec}
\usepackage{fancyhdr}
\usepackage{tcolorbox}
\usepackage{longtable}
\usepackage{array}
\usepackage{booktabs}
\usepackage{hyperref}
\usepackage{microtype}
\usepackage{amssymb}
\tcbuselibrary{skins,breakable,listings}

% --- Colours ---
\definecolor{rocketred}{HTML}{C8102E}
\definecolor{rocketblue}{HTML}{14375F}
\definecolor{warnyellow}{HTML}{F5C518}
\definecolor{gogreen}{HTML}{2E7D32}
\definecolor{nogored}{HTML}{B71C1C}
\definecolor{softgrey}{HTML}{F5F5F5}
\definecolor{linegrey}{HTML}{888888}

\hypersetup{
  colorlinks=true,
  linkcolor=rocketblue,
  urlcolor=rocketblue,
  pdftitle={Pre-Launch Procedures},
  pdfauthor={MPR Altitude Logger Team}
}

% --- Chapter: completely suppressed. We use \rsection (defined later)
%     to render the visible title page, bump the chapter counter, and
%     write the TOC entry — without invoking \chapter's heading machinery.
\makeatletter
\renewcommand{\@makechapterhead}[1]{}
\renewcommand{\@makeschapterhead}[1]{}
\makeatother

% Section / subsection styling
\titleformat{\section}
  {\normalfont\sffamily\Large\bfseries\color{rocketblue}}{\thesection}{0.8em}{}
\titleformat{\subsection}
  {\normalfont\sffamily\large\bfseries\color{rocketblue!85!black}}{\thesubsection}{0.6em}{}

% Headers / footers
\pagestyle{fancy}
\fancyhf{}
\fancyhead[L]{\sffamily\small\textcolor{linegrey}{Pre-Launch Procedures}}
\fancyhead[R]{\sffamily\small\textcolor{linegrey}{\leftmark}}
\fancyfoot[C]{\sffamily\small\thepage}
\renewcommand{\headrulewidth}{0.4pt}
\renewcommand{\footrulewidth}{0pt}

\fancypagestyle{plain}{%
  \fancyhf{}%
  \fancyfoot[C]{\sffamily\small\thepage}%
  \renewcommand{\headrulewidth}{0pt}%
}

% --- Custom callout boxes ---
% NOTE: warnbox and gobox are intentionally NOT breakable so they stay
% on a single page. If they don't fit, they push to the next page.
\newtcolorbox{warnbox}{
  enhanced,
  colback=warnyellow!15, colframe=warnyellow!85!black,
  boxrule=0.8pt, arc=2pt,
  left=8pt, right=8pt, top=4pt, bottom=4pt,
  before skip=10pt, after skip=10pt,
  title={\textbf{WARNING}}, fonttitle=\sffamily\bfseries,
  coltitle=black, colbacktitle=warnyellow,
  attach boxed title to top left={xshift=8pt, yshift=-2pt},
  boxed title style={colback=warnyellow, colframe=warnyellow!85!black, sharp corners=northwest, sharp corners=northeast, arc=2pt}
}

\newtcolorbox{gobox}{
  enhanced,
  colback=gogreen!8, colframe=gogreen,
  boxrule=0.8pt, arc=2pt,
  left=8pt, right=8pt, top=4pt, bottom=4pt,
  before skip=10pt, after skip=10pt,
  title={\textbf{GO / NO-GO}}, fonttitle=\sffamily\bfseries,
  coltitle=white, colbacktitle=gogreen,
  attach boxed title to top left={xshift=8pt, yshift=-2pt},
  boxed title style={colback=gogreen, sharp corners=northwest, sharp corners=northeast, arc=2pt}
}

\newtcolorbox{infobox}{
  enhanced, breakable,
  colback=softgrey, colframe=linegrey,
  boxrule=0.4pt, arc=2pt,
  left=8pt, right=8pt, top=4pt, bottom=4pt
}

% --- Code block (verbatim shell commands) ---
\newtcblisting{codebox}{
  enhanced, breakable,
  listing only,
  colback=black!92, colframe=black!92,
  coltext=white,
  arc=2pt, boxrule=0pt,
  left=10pt, right=10pt, top=4pt, bottom=4pt,
  listing options={
    basicstyle=\ttfamily\small\color{white},
    breaklines=true,
    breakatwhitespace=false,
    columns=fullflexible,
    keepspaces=true,
    showstringspaces=false
  }
}

% --- Inline tweaks ---
% Checkbox bullet for action items (printable, can be ticked by hand)
\newcommand{\cbox}{\raisebox{-0.5pt}{\large$\square$}}
\setlist[itemize,1]{label=\cbox, leftmargin=1.8em, itemsep=3pt, topsep=3pt}
\setlist[itemize,2]{label=\cbox, leftmargin=1.6em, itemsep=2pt, topsep=2pt}
\setlist[enumerate]{leftmargin=1.8em, itemsep=3pt, topsep=3pt}

% Plain-bullet itemize for non-actionable lists (LED reference, abbreviations)
\newlist{plainitem}{itemize}{2}
\setlist[plainitem]{label=\textbullet, leftmargin=1.4em, itemsep=2pt, topsep=2pt}

\newcommand{\done}[1]{\bigskip\begin{center}\sffamily\bfseries\Large\textcolor{rocketred}{#1}\end{center}\bigskip}

% --- Document ---
\begin{document}

% ===== Title page =====
\begin{titlepage}
  \centering
  \vspace*{2cm}
  {\sffamily\bfseries\color{rocketblue}\Huge Pre-Launch Procedures\par}
  \vspace{0.5cm}
  {\color{linegrey}\rule{0.6\textwidth}{1pt}\par}
  \vspace{0.8cm}
  {\sffamily\Large Hobby Rocket Launch\par}
  \vspace{2cm}

  \begin{tabular}{rl}
    \sffamily\bfseries Vehicle: & Hobby rocket \\[4pt]
    \sffamily\bfseries Avionics firmware: & v2.0.0 (MPR Altitude Logger) \\[4pt]
    \sffamily\bfseries Document version: & 1.0 \\[4pt]
    \sffamily\bfseries Date: & \today \\
  \end{tabular}

  \vfill

  \begin{infobox}
    \sffamily\small
    \textbf{Conventions} \\[2pt]
    \textbf{T-N} = N days before launch day. T-0 = launch day at the field. \\
    \textbf{GO / NO-GO} blocks at the end of each phase are mandatory checkpoints --- if any condition fails, do not proceed without the relevant lead's sign-off. \\
    All \texttt{pnpm} commands run from \texttt{\textasciitilde/Documents/mpr-altitude-logger} on Dash's laptop.
  \end{infobox}

  \vspace{1cm}
\end{titlepage}

% ===== Front matter =====
\thispagestyle{plain}
\section*{Abbreviations}
\begin{plainitem}
  \item \textbf{ORK} --- OpenRocket simulation file
  \item \textbf{lwr} --- lower (bulkhead)
  \item \textbf{AGL} --- altitude above ground level
  \item \textbf{WDT} --- watchdog timer (reboot if firmware hangs)
  \item \textbf{SD} --- SD card (microSD on the avionics PCB)
  \item \textbf{TUI} --- text-based UI on the laptop (preflight / postflight tools)
  \item \textbf{PCB} --- printed circuit board (the avionics flight computer)
\end{plainitem}

\tableofcontents
\clearpage

% ===== Section title page macro =====
% \rsection{Title}{Subtitle} — renders a full-page section title,
% bumps the chapter counter, and adds a TOC entry.
\newcommand{\rsection}[2]{%
  \refstepcounter{chapter}%
  \setcounter{section}{0}%
  \addcontentsline{toc}{chapter}{\protect\numberline{\thechapter}#1}%
  \clearpage
  \thispagestyle{empty}%
  \markboth{#1}{}%
  \vspace*{\fill}
  \begin{center}
    {\sffamily\color{linegrey}\Large Section \thechapter}\\[6pt]
    {\color{rocketred}\rule{0.4\textwidth}{1.5pt}}\\[20pt]
    {\sffamily\bfseries\color{rocketblue}\Huge #1}\\[12pt]
    {\sffamily\large\color{linegrey}\itshape #2}
  \end{center}
  \vspace*{\fill}
  \clearpage
}
% Legacy alias so existing \chapter+\sectiontitlepage pairs still work
\newcommand{\sectiontitlepage}[2]{}

% --- Section page breaks ---
% \section{} (T-1, T-0, Post-flight, Field Recovery, Comms, etc.) starts on a new page.
% \proc{}  is a procedure heading that flows inline (no forced page break).
\let\rsoldsection\section
\renewcommand{\section}[1]{\clearpage\rsoldsection{#1}}

\newcommand{\proc}[1]{\subsection*{#1}}

% =====================================================================
\rsection{Aerostructures}{Airframe build, ballast, integration}

\section{T-1}

\proc{Check wind predictions for the day}
\begin{itemize}
  \item Pull forecast for the launch site, hour-by-hour for the launch window
  \item Conduct additional ORK simulations at predicted windspeed $\pm$ 20\%
  \item Calculate appropriate ballast for windspeed at predicted $\pm$ 50\%
  \item Record the ballast mass for each scenario in the team sheet
\end{itemize}

\section{T-0 (Launch Day)}

\proc{Install eyebolt in lower bulkhead}
\begin{itemize}
  \item Drill as close to the centre as possible
  \item Use the pre-marked drill jig if available
  \item Confirm eyebolt threads engage cleanly before epoxy step
\end{itemize}

\proc{Motor integration}
\begin{itemize}
  \item Friction-fit the motor into the airframe
  \item Use the motor retention clip to retain the motor
  \item Confirm clip is fully seated --- no visible gap
\end{itemize}

\proc{Measure wind}
\begin{itemize}
  \item Use the handheld anemometer (the ``fan-lookin' thing'')
  \item Record three readings 30 seconds apart at pad height
  \item If wind exceeds the launch limit from your ORK sims, hold
\end{itemize}

\proc{Add ballast to the tip of the nose cone}
\begin{itemize}
  \item Mix up some 5-minute epoxy
  \item Log weight of mixed epoxy
  \item Measure split shot weight: $W_{ss} = W_{ballast} - W_{epoxy}$
  \item Place split shot sinkers at the tip of the nose cone, then add the epoxy
  \item Hold the nosecone nose-down until the epoxy cures ($\sim$5 min, verify with toothpick)
\end{itemize}

\proc{Add battery (coordinate with Avionics --- battery only goes in AFTER avionics bench check passes)}
\begin{itemize}
  \item Connect battery to the flight computer
  \item Battery life on fresh 9V alkaline: $\sim$4+ hours idle. Install within 1 hour of launch slot. Set a phone timer.
  \item Lower battery into the slot in the nose cone
  \item Secure battery: electrical tape preferred. Hot glue is a fallback only --- accept that it may shift under boost loads.
\end{itemize}

\proc{Install flight computer}
\begin{itemize}
  \item Lower the connected PCB into the nosecone
  \item Use 3$\times$ M3 magnetic screws + screwdriver to screw the PCB in place
  \item DO NOT overtighten --- PLA threads strip easily
  \item Verify slow-blink LED is still visible through the nosecone vent / pressure hole before sealing
  \item Note: the lower bulkhead glue (next step) seals the bay. Post-flight recovery requires cutting the nosecone open --- the screws are only accessible after that cut.
\end{itemize}

\proc{Install lower bulkhead}
\begin{itemize}
  \item Mix up 5-minute epoxy
  \item Lather the sides and rear of the bulkhead with epoxy
  \item Using the eyebolt, hold the bulkhead in place until epoxy cures
  \item Tie the shock cord to the eyebolt
\end{itemize}

\done{NOSECONE DONE!}

\proc{Integrate nosecone with main airframe}

\proc{Drill pressure hole}
\begin{itemize}
  \item Take care not to drill into the avionics bay or shock cord
  \item Hole diameter per design spec --- confirm before drilling
\end{itemize}

\begin{gobox}
\textbf{GO:} ballast logged, motor retained, bulkhead cured, pressure hole drilled, avionics LED slow-blinking through pressure hole. \\[4pt]
\textbf{NO-GO:} any epoxy uncured, motor not fully seated, avionics LED solid or absent.
\end{gobox}

\done{STRUCTURES DONE}

% =====================================================================
\rsection{Avionics}{Firmware, preflight, GO/NO-GO, recovery cut-out}

\section{T-1}

\proc{Firmware deploy (Dash's laptop, board on USB)}

\begin{codebox}
cd ~/Documents/mpr-altitude-logger
git status            # clean tree, or know what's dirty
pnpm version          # must print: Firmware version: 2.0.0
pnpm deploy:pico      # wait for VERIFY_OK at the bottom
\end{codebox}

\begin{itemize}
  \item If any module reports \texttt{FAIL}, STOP --- fix on the bench, do not bring a half-deployed board to the field
\end{itemize}

\proc{Bench preflight (smoke test on home network/desk before packing)}
\begin{itemize}
  \item \texttt{pnpm preflight}
  \item Confirm all rails green, I2C 0x77 present, BMP180 chip ID OK, RAM free $>$ 50 KB
  \item Press \texttt{[Q]} to quit cleanly
\end{itemize}

\proc{10-minute logging stress test (T-1 only --- too slow for launch day)}

The board flies for less than a minute, but a degraded SD or marginal timing only shows up over many minutes. Run this once the night before.

\begin{itemize}
  \item \texttt{pnpm preflight} $\to$ \texttt{[R]} recalibrate $\to$ \texttt{[B]} boot to PAD
  \item Disconnect USB. Power the board from a 9V (or leave on USB if you trust the host won't sleep) for at least 10 minutes
  \item LED should be slow-blink (PAD) the whole time. If it ever goes solid, abort the test and investigate.
  \item After 10 min, power off, pull SD, insert into laptop, then:
\end{itemize}

\begin{codebox}
pnpm decode -- "/Volumes/AVIONICS/AVIONICS_flight_NNN.bin" --plot
\end{codebox}

\begin{itemize}
  \item Verify, at 50 Hz sample rate:
  \begin{itemize}
    \item Frame count $\approx$ 30,000 $\pm$ 1\% over 600 s (i.e. 29,700--30,300 frames)
    \item Mean sample interval = 20 ms $\pm$ 1 ms
    \item No timing gap $>$ 50 ms (one missed frame is suspicious; multiple = bad)
    \item No \texttt{[CRASH REBOOT]} line in the boot log on next power-up
    \item File grew monotonically (no truncation), no \texttt{\textbackslash xAA\textbackslash x55} resync gaps in the middle
  \end{itemize}
  \item If frame count is short, intervals are off, or there are gaps, do NOT fly --- swap SD card and re-run. If a fresh SD also fails, investigate firmware before bringing the board to the field.
\end{itemize}

\proc{SD card prep}
\begin{itemize}
  \item Confirm SD seated in the avionics PCB
  \item FAT32 formatted, $>$ 50 MB free (check the existing \texttt{flights/} folder isn't full)
  \item Spare SD card formatted FAT32 and packed
\end{itemize}

\proc{Battery prep}
\begin{itemize}
  \item 2$\times$ fresh 9V alkaline (one flight, one spare)
  \item Unloaded multimeter check --- reject anything under 9.0 V
  \item Tape over terminals until install
\end{itemize}

\section{T-0 (Launch Day)}

\proc{Bench preflight at the field (do this BEFORE Aerostructures starts nosecone integration)}
\begin{itemize}
  \item Connect Pico to laptop via USB. Wait for \texttt{/dev/cu.usbmodem*} to appear.
\end{itemize}

\begin{codebox}
cd ~/Documents/mpr-altitude-logger
pnpm preflight
\end{codebox}

\begin{itemize}
  \item Verify EVERY line is green:
  \begin{itemize}
    \item I2C scan finds device 0x77
    \item BMP180 chip ID OK
    \item Pressure within $\pm$5 hPa of weather report (sanity-check the sensor)
    \item 3V3 rail in 3.0--3.6 V
    \item 5V rail in 4.5--5.5 V
    \item 9V rail in 8.0--10.0 V --- if low, swap battery now
    \item SD mounted, free space $>$ 50 MB
    \item Frame size = 40 bytes (v3 log format)
  \end{itemize}
  \item If anything is red, run \texttt{pnpm pico:diag} for raw live diagnostics
\end{itemize}

\proc{Calibrate and arm}
\begin{itemize}
  \item Press \texttt{[R]} to recalibrate ground pressure at pad elevation
  \item Watch the live altitude trace for 10--15 seconds standing still
  \item Filtered altitude must hold within $\pm$1 m and not drift. If it drifts, recalibrate. If it still drifts, abort to NO-GO.
  \item Press \texttt{[B]} to boot into flight mode
  \item TUI exits cleanly. LED transitions to slow blink (1 s on / 1 s off) = PAD state, armed, logging-ready
  \item Disconnect USB. Do NOT pull power yet --- flight battery only from here on
\end{itemize}

\proc{Hand off to Aerostructures}
\begin{itemize}
  \item Coordinate the ``Add battery'' $\to$ ``Install flight computer'' $\to$ ``Install lower bulkhead'' sequence
  \item Battery should be installed within 1 hour of expected launch slot --- phone timer mandatory
  \item After battery is connected and PCB installed, confirm slow-blink LED is visible through the nosecone vent before Aerostructures closes the bulkhead
\end{itemize}

\proc{Final pad check (just before vertical)}
\begin{itemize}
  \item Visually confirm slow-blink LED through the pressure hole / vent
  \item Slow blink only = GO. Anything else = NO-GO
  \item No laptop interaction at the pad. The board self-detects launch (alt $>$ 15 m AND velocity $>$ 10 m/s, sustained 0.5 s)
\end{itemize}

\proc{LED reference (memorise this)}
\begin{plainitem}
  \item Fast blink (250 ms) --- booting / preflight in progress
  \item Solid ON --- error
  \item \textbf{Slow blink (1 s) --- PAD, ready, safe to fly}
  \item Fast blink (50 ms) --- BOOST detected
  \item Medium blink --- COAST / DROGUE / MAIN descent
  \item Double flash --- APOGEE
  \item Triple flash --- LANDED, data safe
\end{plainitem}

\section{Post-flight (after recovery)}

The avionics PCB is held by 3$\times$ M3 screws, but the lower bulkhead is epoxied in place --- so the bay is sealed and the screws cannot be reached without cutting the nosecone open. Power is still live until the battery is unplugged --- treat the board as energised.

\proc{Cut the nosecone}
Choose ONE cut location:
\begin{itemize}
  \item \textbf{Option A} --- circumferential cut \emph{just below the avionics} (between PCB and the lower bulkhead). Exposes the battery side first.
  \item \textbf{Option B} --- circumferential cut \emph{just above the avionics} (between PCB and the nose tip / ballast). Exposes the PCB top side first.
\end{itemize}

\begin{itemize}
  \item Mark the cut line with a pencil all the way around before cutting
  \item Use a hand saw, rotary tool, or sharp knife --- slow and steady, multiple shallow passes
  \item Pause and inspect every few mm. Pull debris out as you go.
\end{itemize}

\begin{warnbox}
Stop cutting the moment you feel the wall give. Do NOT push the blade past the wall thickness --- the board, battery, or wiring is directly behind it.
\end{warnbox}

\proc{Disconnect the battery FIRST}
\begin{itemize}
  \item Once the cut is open, locate the battery clip and disconnect it before doing anything else
  \item LED should go dark --- confirms power is off
  \item Tape over the 9V terminals immediately
\end{itemize}

\proc{Free the PCB and SD}
\begin{itemize}
  \item Brush or blow swarf out of the bay before unscrewing
  \item Reach through the cut and unscrew the 3$\times$ M3 PCB screws
  \item Catch each screw as it comes out --- if one drops into the bay it can roll under the PCB and short it. Account for all 3.
  \item Lift the PCB out cleanly, supporting it by the edges (not by components)
  \item Pull the SD card out cleanly
\end{itemize}

\begin{warnbox}
Do NOT let a screw, screwdriver tip, blade, or piece of plastic swarf touch the board while power is live or after. Debris bridging pads can short the board and brick it.
\end{warnbox}

\proc{Decode the flight}
\begin{itemize}
  \item Insert the SD card into laptop --- auto-mounts as \texttt{/Volumes/AVIONICS/}
  \item \texttt{pnpm postflight} --- TUI auto-detects SD, lists flights newest first, decodes selected flight
  \item Save the \texttt{.bin}, \texttt{.csv}, and \texttt{\_report.txt} to the \texttt{flights/} folder
  \item If summary looks anomalous: \texttt{pnpm decode -- "flights/AVIONICS\_flight\_NNN.bin" --plot} to inspect raw frames
  \item If a \texttt{[CRASH REBOOT]} line appears in the boot log, save the file and flag it before the next flight --- do not reuse the SD until investigated
\end{itemize}

\proc{Multi-flight day --- between flights}
\begin{itemize}
  \item Disconnect avionics power, eject and reseat SD, swap or top up 9V
  \item USB $\to$ laptop $\to$ \texttt{pnpm preflight} $\to$ \texttt{[R]} recalibrate $\to$ \texttt{[B]} boot
  \item Disconnect $\to$ into airframe
  \item Each flight gets an auto-incremented filename (\texttt{AVIONICS\_flight\_NNN.bin}) --- no overwrite risk
\end{itemize}

\begin{gobox}
\textbf{GO:} all preflight rails green, ground pressure calibrated, slow-blink LED on PAD, SD mounted with $>$ 50 MB free. \\[4pt]
\textbf{NO-GO} (any of these = scratch the flight):
\begin{itemize}
  \item Solid LED at any point after preflight
  \item 9V rail $<$ 8.5 V on TUI
  \item Altitude drift $>$ 2 m on bench
  \item SD won't mount or $<$ 10 MB free
  \item Frame timing $>$ 30 ms in pico\_diag
  \item Stuck fast-blink LED that won't settle to slow blink within 30 s of boot
\end{itemize}
\end{gobox}

\proc{Emergency: avionics won't boot at the pad}
\begin{enumerate}
  \item Solid LED $\to$ known-bad state. No spare PCB --- scratch the flight.
  \item No LED at all $\to$ power. Check 9V, switch, connector.
  \item Stuck fast-blink $\to$ power cycle once. If it repeats, swap SD card first (cheapest, fastest fix).
\end{enumerate}

\textit{When in doubt, scratch. A scratched flight is recoverable; a corrupt log is not.}

\done{AVIONICS DONE!}

% =====================================================================
\rsection{Recovery}{Chute pack, field recovery, avionics handling}

\begin{infobox}
\textbf{Deployment:} chute is deployed by the motor's delay charge. There are no separate pyro charges, e-matches, or electronic deployment. The motor's delay grain ejects the chute via the forward closure when the propellant burns through.
\end{infobox}

\section{T-1}

\proc{Chute inspection}
\begin{itemize}
  \item Unfold chute fully, lay flat
  \item Check canopy for tears, scorch marks, melted shroud lines
  \item Check shock cord for fraying, knots holding firm
  \item Check swivels and quick links for damage
  \item Repack chute loosely (final pack happens T-0)
\end{itemize}

\section{T-0 (Launch Day)}

\proc{Chute pack and install}
\begin{itemize}
  \item Fold chute per team's standard fold (Z-fold or roll, whatever you train on)
  \item Wrap in nomex or chute protector to shield from motor exhaust gases
  \item Connect to shock cord, shock cord to eyebolt on lower bulkhead
  \item Pack into airframe, leaving room above for ejection
  \item Confirm motor delay matches predicted apogee time from ORK sim --- if delay is too short the chute deploys under high speed and shreds; too long and the rocket lawn-darts
\end{itemize}

\begin{warnbox}
The motor delay must match the predicted coast time. Verify the motor designation includes the right delay (e.g.\ \texttt{H100-7} = 7-second delay). A wrong delay = destroyed rocket. Check twice before motor install.
\end{warnbox}

\begin{gobox}
\textbf{GO:} chute inspected, packed, shock cord connected, motor delay matches sim ($\pm$1 s). \\[4pt]
\textbf{NO-GO:} damaged canopy, frayed shock cord, motor delay mismatch.
\end{gobox}

\section{Field Recovery (post-launch)}

\proc{Watch and track}
\begin{itemize}
  \item Eyes on the rocket from launch to landing --- assign a designated tracker
  \item Note apogee location and chute deployment (visible puff)
  \item Track descent direction; note landmarks near landing zone
  \item Wait for Range Director's clearance before walking onto the field
\end{itemize}

\proc{Walk-out}
\begin{itemize}
  \item Recovery team: 2+ people, hi-vis vests, phone with GPS, radio if available
  \item Don't cross active landing zones for other rockets until cleared
  \item Photograph the landing site before touching --- captures orientation and any damage
\end{itemize}

\proc{Approach the rocket}
\begin{itemize}
  \item Spent motor case may still be hot --- give it a few minutes if you arrive fast
  \item Inspect for damage: bent fins, cracked airframe, tangled chute, missing nosecone
  \item Check chute condition for post-flight notes
\end{itemize}

\proc{Avionics handling in the field}

The avionics is still powered. The board's last act is to detect LANDED and flush the SD --- the LED tells you whether that worked.

\begin{itemize}
  \item Look at the LED through the nosecone vent / pressure hole:
  \begin{itemize}
    \item \textbf{Triple flash} = LANDED, data flushed, safe to transport. Best case.
    \item \textbf{Slow blink} = still in PAD --- launch was never detected, log will be short or empty
    \item \textbf{Other pattern} = stuck in a mid-flight state, may not have flushed cleanly
  \end{itemize}
  \item DO NOT disconnect the battery in the field unless you have to. Let the board sit at LANDED so any pending writes finish.
  \item If the LED is in an odd state, photograph it before moving --- useful for post-flight diagnosis
\end{itemize}

\begin{warnbox}
Vibration during SD writes can corrupt the last few frames. If the board is still mid-flush (slow LED, not yet triple flash), avoid jostling the rocket. Walking it back is fine; throwing it in the back of a car is not.
\end{warnbox}

\proc{Transport back to bench}
\begin{itemize}
  \item Carry the rocket supported, not by the nosecone (cumulative shock isn't great for the SD)
  \item Battery stays connected. If transport is $>$ 30 min and the LED is showing triple flash (LANDED), you can power off if needed --- the data is safe.
  \item Bag the rocket if it's wet, muddy, or going through dust
  \item Once at the bench, hand off to Avionics for the post-flight cut-out (see Avionics section)
\end{itemize}

\proc{If the rocket is lost / unrecoverable}
\begin{itemize}
  \item Mark the last seen position on a map (phone GPS pin)
  \item Search in expanding circles from predicted landing zone
  \item Ask other teams if they spotted it
  \item If unrecoverable: the SD card and battery are in the rocket, and the data is gone. Note for post-mortem.
\end{itemize}

\done{RECOVERY DONE}

% =====================================================================
\rsection{Propulsions}{Motor, igniter, paperwork}

\begin{infobox}
\textbf{TO FILL IN --- owner: \rule{4cm}{0.4pt}}
\end{infobox}

Suggested structure:
\begin{plainitem}
  \item \textbf{T-1:} motor inspection, igniter inventory, certification paperwork
  \item \textbf{T-0:} motor install (coordinate with Aerostructures), igniter install, continuity check, GO/NO-GO criteria
  \item \textbf{Post-flight:} motor case retrieval, post-fire inspection
\end{plainitem}

% =====================================================================
\rsection{Master Timeline (T-0)}{Sequencing and comms}

Approximate, adjust to launch slot:

\renewcommand{\arraystretch}{1.3}
\begin{longtable}{@{}p{3cm} p{4.5cm} p{8cm}@{}}
\toprule
\textbf{Time} & \textbf{Owner} & \textbf{Action} \\
\midrule
\endhead
Slot - 3:00 h & All & Arrive, unpack, set up bench \\
Slot - 2:30 h & Aerostructures & Eyebolt, motor friction fit, ballast prep \\
Slot - 2:00 h & Avionics & Bench preflight, calibrate, boot to PAD \\
Slot - 1:30 h & Avionics + Aerostructures & Battery install, PCB install in nosecone \\
Slot - 1:15 h & Aerostructures & Lower bulkhead epoxy, nosecone integration \\
Slot - 0:45 h & Recovery + Propulsions & Charges, igniter, final continuity \\
Slot - 0:30 h & All & Walk to pad, vertical, final visual checks \\
Slot - 0:10 h & Range & Range crew takeover \\
\textbf{Slot 0:00} & --- & \textbf{LAUNCH} \\
Slot + 0:30 h & All & Recovery walk \\
Slot + 1:00 h & Avionics & Cut nosecone, disconnect battery, pull SD, postflight decode \\
\bottomrule
\end{longtable}

\section{Comms}
\begin{itemize}
  \item \textbf{Launch Director} calls GO/NO-GO at each phase. Anyone can call NO-GO at any time for any reason.
  \item Subsystem leads must verbally confirm GO at handoff points (Avionics $\to$ Aerostructures battery install, etc.)
\end{itemize}

% =====================================================================
\rsection{Pack List}{What goes in the box}

\section{Structures}
\begin{itemize}
  \item Spare nosecone
  \item 5-minute epoxy (lower bulkhead)
  \item 3$\times$ M3 magnetic screws (avionics PCB) + spares
  \item Magnetic screwdriver (M3 bit) --- PCB install + post-flight removal
  \item Split shot sinkers
  \item Eye bolt
  \item Drill + drill bits (eyebolt + pressure hole)
  \item Toothpicks (epoxy cure check)
  \item Electrical tape, hot glue gun + sticks (fallback)
  \item Hand saw / Dremel / sharp knife (for post-flight nosecone cut)
  \item Bench vice or clamp (to hold nosecone steady during cut)
  \item Safety glasses
\end{itemize}

\section{Avionics}
\begin{itemize}
  \item Laptop + charger (Dash)
  \item USB-A to micro-USB cable --- DATA capable, not charge-only
  \item 2$\times$ fresh 9V batteries (alkaline)
  \item Multimeter
  \item Spare SD card, formatted FAT32
  \item USB SD card reader
  \item Printout of this doc (no field wifi guaranteed)
\end{itemize}

\section{Recovery}
\begin{itemize}
  \item Parachute (main + drogue if applicable)
  \item Shock cord, swivels
  \item Ejection charges (per manifest)
  \item Igniter / e-match wire
  \item Continuity tester
  \item Nomex / chute protector
  \item Spare quick links
\end{itemize}

\section{Propulsions}
\begin{itemize}
  \item Motor (per manifest)
  \item Spare igniter(s)
  \item Motor retention clip
  \item Cert paperwork
  \item Igniter wire
  \item Continuity tester (shared with Recovery OK)
\end{itemize}

\section{Shared}
\begin{itemize}
  \item First aid kit
  \item Hi-vis vests
  \item Sunscreen, water, snacks
  \item Folding table, chairs
  \item Power bank for laptop
\end{itemize}

\end{document}
